本节从独立的无反射离散数据出发, 不预设待构造三场已经满足 HMB 方程.
先求一般有限维系统, 再对指定纯虚轨道证明全局可逆性.
本节的常数均取初始时刻, 时间因素由式 \eqref{eq:C-timephase} 统一给出.

\subsection{多轨道无反射问题的条件性求解}
这里的 $N$ 按完整简单四点轨道的数目计数. 以下在无反射条件下求出矩阵解及相应场的重构公式.

当 $\rho=\widetilde\rho=0$ 时, 跳跃矩阵为 $I$. 在非极点轮廓处取可去条件. 在分支点 $z_b=\pm\ii q_0$ 处, 矩阵 $M$ 正则, 而其行列式 $\gamma(z)=1+\dfrac{q_0^2}{z^2}$ 具有一阶零点, 即 $\gamma(z_b)=0$、$\gamma'(z_b)\ne0$.

固定 $N$ 个两两不交的完整简单四点轨道. 按 \eqref{eq:discrete-indexing} 选取外部上半平面代表 $\zeta_1,\ldots,\zeta_N$, 并置
\begin{equation}\label{eq:N-orbit-data}
\begin{aligned}
m&=2N,\qquad \zeta_{N+n}=-\dfrac{q_0^2}{\conj{\zeta_n}},\qquad n=1,\ldots,N,\\
c_{N+n}(0)&=-\dfrac{q_0^2}{(\conj{\zeta_n})^2}\conj{c_n(0)},\qquad
C_j=c_j(0)e^{2\ii\theta(\zeta_j)},\qquad j=1,\ldots,m.
\end{aligned}
\end{equation}
所有 $4N$ 个谱点互异且位于 $\Sigma$ 外, 各 $c_n(0)\ne0$. 同相背景要求
\begin{equation}\label{eq:N-orbit-phase}
\prod\limits_{j=1}^{m}\dfrac{\zeta_j}{\conj{\zeta_j}}=1.
\end{equation}
下文保留全部 $m$ 个 $D_+$ 点及其共轭点, 令 $\mathbf v_j=M_1(\zeta_j)$、$\mathbf u_j=M_2(\conj{\zeta_j})$.

\paragraph{留数系统与分块消元.}\label{prop:finite}
除 $0$ 与指定离散极点外没有有限奇点的无反射矩阵解具有表示
\begin{equation}\label{eq:partialfraction}
M(z)=I-\dfrac{\ii q_0}{z}\sigma_1+\sum\limits_{j=1}^{m}
\left(\dfrac{C_j\mathbf v_je_2^T}{z-\zeta_j}
-\dfrac{\conj{C_j}\mathbf u_je_1^T}{z-\conj{\zeta_j}}\right),
\end{equation}
其正则列满足
\begin{align}
\mathbf v_j&=e_1-\dfrac{\ii q_0}{\zeta_j}e_2
-\sum\limits_{k=1}^{m}\dfrac{\conj{C_k}\mathbf u_k}{\zeta_j-\conj{\zeta_k}},\label{eq:finite1}\\
\mathbf u_j&=e_2-\dfrac{\ii q_0}{\conj{\zeta_j}}e_1
+\sum\limits_{k=1}^{m}\dfrac{C_k\mathbf v_k}{\conj{\zeta_j}-\zeta_k}.\label{eq:finite2}
\end{align}
反之, 任一系统解通过 \eqref{eq:partialfraction} 给出具有指定留数和两处归一化的矩阵. 记
\[
A_{jk}=\dfrac{\conj{C_k}}{\zeta_j-\conj{\zeta_k}},\qquad
B_{jk}=\dfrac{C_k}{\conj{\zeta_j}-\zeta_k},\qquad j,k=1,\ldots,m,
\]
并记 $f_j=e_1-\dfrac{\ii q_0}{\zeta_j}e_2$、$g_j=e_2-\dfrac{\ii q_0}{\conj{\zeta_j}}e_1$. 将 $\mathbf v_j^T,\mathbf u_j^T,f_j^T,g_j^T$ 分别作为 $m\times2$ 矩阵 $\boldsymbol v,\boldsymbol u,\boldsymbol f,\boldsymbol g$ 的第 $j$ 行. 则留数系统为
\begin{equation}\label{eq:N-block-system}
K_N\begin{pmatrix}\boldsymbol v\\\boldsymbol u\end{pmatrix}
=\begin{pmatrix}\boldsymbol f\\\boldsymbol g\end{pmatrix},\qquad
K_N=\begin{pmatrix}I_m&A\\-B&I_m\end{pmatrix}.
\end{equation}
置 $L_N=I_m+AB$. 若 $\det L_N\ne0$, 则系统的唯一解为
\begin{equation}\label{eq:finite-schur}
\boldsymbol v=L_N^{-1}(\boldsymbol f-A\boldsymbol g),\qquad
\boldsymbol u=\boldsymbol g+B\boldsymbol v.
\end{equation}
相应光场为
\begin{equation}\label{eq:qfinite}
q=q_0+\ii\sum\limits_{j=1}^{m}C_j(\mathbf v_j)_1.
\end{equation}

上述关系由以下计算得到.

减去 $z=0$ 处主部及全部离散主部后, 剩余矩阵在有限平面解析, 且在无穷远趋于 $I$. Liouville 定理给出 \eqref{eq:partialfraction}, 这也正是 \eqref{eq:general-cauchy} 中 $G=0$ 的情形. 在第一列取 $z=\zeta_j$, 在第二列取 $z=\conj{\zeta_j}$, 即得 \eqref{eq:finite1}--\eqref{eq:finite2}. 反向代入时, 这些取值等式保证留数系数恰为对应正则列. 写成
\[
\boldsymbol v=\boldsymbol f-A\boldsymbol u,\qquad
\boldsymbol u=\boldsymbol g+B\boldsymbol v,
\]
对系数矩阵作分块消元,
\[
\begin{pmatrix}I_m&-A\\0&I_m\end{pmatrix}
\begin{pmatrix}I_m&A\\-B&I_m\end{pmatrix}
=\begin{pmatrix}L_N&0\\-B&I_m\end{pmatrix},
\qquad \det K_N=\det L_N.
\]
当此行列式非零时, 第一行解出 $\boldsymbol v$, 第二行给出 $\boldsymbol u$, 得到 \eqref{eq:finite-schur}. 最后读取 \eqref{eq:partialfraction} 的 $(1,2)$ 元的 $z^{-1}$ 系数, 得到 \eqref{eq:qfinite}.


\begin{theorem}[$N$ 轨道矩阵解与行列式重构]\label{thm:N-orbit-RHP}
设数据满足 \eqref{eq:N-orbit-data}--\eqref{eq:N-orbit-phase}. 固定实点 $(x,t)$ 并假设留数矩阵满足 $\det L_N(x,t)\ne0$. 则式 \eqref{eq:finite-schur} 和 \eqref{eq:partialfraction} 给出无反射 RHP 的唯一解. 它满足 \eqref{eq:Msym} 及 $\det M=\gamma$, 在分支点解析且秩为一. 令
\begin{equation}\label{eq:N-rhs}
\begin{aligned}
&\boldsymbol C=(C_1,\ldots,C_m)^T,\qquad
\boldsymbol r=(\boldsymbol f-A\boldsymbol g)e_1,\\
&r_j=1+\ii q_0\sum\limits_{k=1}^{m}
\dfrac{\conj{C_k}}{(\zeta_j-\conj{\zeta_k})\conj{\zeta_k}},\qquad j=1,\ldots,m,
\end{aligned}
\end{equation}
则规范光场有显式矩阵式与行列式表示
\begin{equation}\label{eq:N-determinant-q}
q=q_0+\ii\boldsymbol C^T L_N^{-1}\boldsymbol r
=\dfrac{\det\begin{pmatrix}L_N&\boldsymbol r\\-\ii\boldsymbol C^T&q_0\end{pmatrix}}
{\det L_N}.
\end{equation}
介质矩阵由同一有理矩阵在共振原像处重构:
\begin{equation}\label{eq:N-medium}
R=M(z_a^\pm)\dfrac{\eta_0\lambda(z_a^\pm)}{a}\sigma_3 M(z_a^\pm)^{-1}.
\end{equation}
两个符号给出相同结果; 相应三场按 \eqref{eq:original-fields} 恢复.
\end{theorem}
\begin{proof}
上述留数系统 给出有理矩阵的全部留数及归一化, 并由 $\det L_N\ne0$ 保证唯一性. 共轭变换 $\sigma_2\conj{M(\conj{z})}\sigma_2$ 交换两组极点, 满足同一留数系统和两处归一化. 换叶变换
\[
\widehat M(z)=M(h(z))\left(\dfrac{\ii z}{q_0}\sigma_1\right)^{-1}
\]
也满足这些条件. 具体地, 记同一轨道的配对指标为 $\widehat j$, 即 $\zeta_{\widehat j}=h(\conj{\zeta_j})$. 在 $z=\zeta_j$ 处, $h(z)$ 趋于 $\conj{\zeta_{\widehat j}}$, 由 $h'(\zeta_j)=\dfrac{q_0^2}{\zeta_j^2}$ 和 \eqref{eq:Cpair},
\[
\Res_{z=\zeta_j}\widehat M_2(z)
=\dfrac{\ii q_0\conj{C_{\widehat j}}}{\zeta_jh'(\zeta_j)}\mathbf u_{\widehat j}
=C_j\left(-\dfrac{\ii q_0}{\zeta_j}\mathbf u_{\widehat j}\right)
=C_j\widehat M_1(\zeta_j).
\]
共轭点处同理, 得到带负号的第一列留数条件. 在 $0,\infty$ 处, 换叶交换两处归一化, 仍给出规定主项. 由唯一性得到两项约化 \eqref{eq:Msym}.

在每个离散极点, 奇异列的留数为正则列的标量倍数, 因而 $\det M$ 的留数为零. 又由 $0,\infty$ 的归一化,
\[
\det M(z)=1+\dfrac{d}{z}+\dfrac{q_0^2}{z^2}
\]
对某个标量 $d$ 成立. 换叶约化的行列式给出
\[
1-\dfrac{dz}{q_0^2}+\dfrac{z^2}{q_0^2}
=\det M(h(z))=\dfrac{z^2}{q_0^2}\det M(z)
=\dfrac{z^2}{q_0^2}+\dfrac{dz}{q_0^2}+1,
\]
故 $d=0$, 即 $\det M=\gamma$. 分支点不在部分分式的极点集合中, 所以 $M$ 在该处解析. 若 $M(z_b)$ 为零矩阵, 则行列式至少有二阶零点, 与 $\gamma'(z_b)\ne0$ 矛盾; 因此其秩恰为一.

由 \eqref{eq:finite-schur}, 正则列第一分量排成的列向量为 $L_N^{-1}\boldsymbol r$. 代入 \eqref{eq:qfinite} 得到 \eqref{eq:N-determinant-q} 的第一式. 再对其分子作分块行列式计算,
\[
\det\begin{pmatrix}L_N&\boldsymbol r\\-\ii\boldsymbol C^T&q_0\end{pmatrix}
=\det L_N\left(q_0+\ii\boldsymbol C^T L_N^{-1}\boldsymbol r\right),
\]
得到第二式. 共振原像不在离散谱中, 且 $\gamma(z_a^\pm)>0$, 所以逆矩阵存在. 两个原像互为换叶, $\lambda(h(z))=-\lambda(z)$, 而 $\sigma_1\sigma_3\sigma_1=-\sigma_3$; 代入换叶约化即得 \eqref{eq:N-medium} 与谱叶无关.
\end{proof}

% BEGIN EXPANSION 08-N
\subsubsection*{完整多轨道数据、Schur 消元与行列式公式的逐项计算}
\paragraph{一个轨道为什么贡献两个正则列指标.}
每个外部上半平面的代表 $\zeta_n$ 对应四点
\[
\zeta_n,\quad \conj{\zeta_n},\quad
h(\zeta_n)=-\dfrac{q_0^2}{\zeta_n},\quad
h(\conj{\zeta_n})=-\dfrac{q_0^2}{\conj{\zeta_n}}.
\]
其中第一点与第四点属于 $D_+$, 第二点与第三点属于 $D_-$.
因此 $N$ 个完整轨道共有 $m=2N$ 个 $D_+$ 点,
而不是只取 $N$ 个外部代表就结束求和.
第 $n$ 与第 $N+n$ 个指标分别指第一点和第四点;
它们的共轭点自动补齐另一个解析域.
因为 $\zeta_n\notin\Sigma$, $\zeta_n\ne0$,
且不同轨道不相交, 所有差 $\zeta_j-\conj{\zeta_k}$ 均不为零.
所以 $A,B$ 的矩阵元在给定数据下都有定义.

由 $\theta(h(\conj{\zeta_n}))=-\conj{\theta(\zeta_n)}$,
\[
\begin{aligned}
C_{N+n}
&=-\dfrac{q_0^2}{(\conj{\zeta_n})^2}\conj{c_n(0)}
e^{-2\ii\conj{\theta(\zeta_n)}}\\
&=-\dfrac{q_0^2}{(\conj{\zeta_n})^2}\conj{C_n}.
\end{aligned}
\]
这说明时空指数保持范数常数的配对关系.
同一轨道对相位乘积的贡献逐因子为
\[
\dfrac{\zeta_n}{\conj{\zeta_n}}
\dfrac{\zeta_{N+n}}{\conj{\zeta_{N+n}}}
=\dfrac{\zeta_n}{\conj{\zeta_n}}
\dfrac{-\dfrac{q_0^2}{\conj{\zeta_n}}}{-\dfrac{q_0^2}{\zeta_n}}
=\left(\dfrac{\zeta_n}{\conj{\zeta_n}}\right)^2.
\]
于是 \eqref{eq:N-orbit-phase} 与上一小节的相位和条件一致.
该约束规定两端背景的相位, 不等同于有限维矩阵的可逆性条件.

\paragraph{把有理矩阵的两列与四个元全部写出.}
由部分分式式 \eqref{eq:partialfraction},
\[
\begin{aligned}
M_1(z)&=e_1-\dfrac{\ii q_0}{z}e_2
-\sum\limits_{k=1}^{m}\dfrac{\conj{C_k}\mathbf u_k}{z-\conj{\zeta_k}},\\
M_2(z)&=e_2-\dfrac{\ii q_0}{z}e_1
+\sum\limits_{k=1}^{m}\dfrac{C_k\mathbf v_k}{z-\zeta_k}.
\end{aligned}
\]
因此
\[
\begin{aligned}
M_{11}(z)&=1-\sum\limits_{k=1}^{m}
\dfrac{\conj{C_k}(\mathbf u_k)_1}{z-\conj{\zeta_k}},\\
M_{21}(z)&=-\dfrac{\ii q_0}{z}-\sum\limits_{k=1}^{m}
\dfrac{\conj{C_k}(\mathbf u_k)_2}{z-\conj{\zeta_k}},\\
M_{12}(z)&=-\dfrac{\ii q_0}{z}+\sum\limits_{k=1}^{m}
\dfrac{C_k(\mathbf v_k)_1}{z-\zeta_k},\\
M_{22}(z)&=1+\sum\limits_{k=1}^{m}
\dfrac{C_k(\mathbf v_k)_2}{z-\zeta_k}.
\end{aligned}
\]
第一列在每个 $\zeta_j$ 处正则, 第二列在每个 $\conj{\zeta_j}$ 处正则.
分别取值就得到 \eqref{eq:finite1}--\eqref{eq:finite2}.
反向代入时, 第二列在 $\zeta_j$ 的留数是 $C_j\mathbf v_j$,
而系统保证 $\mathbf v_j=M_1(\zeta_j)$;
第一列的负号也由相同取值核对, 没有额外的未知留数.

\paragraph{为什么 $4N$ 阶系统能缩减为 $2N$ 阶系统.}
$\boldsymbol v,\boldsymbol u$ 各有 $m$ 行, 每行有两个标量分量.
对任意一个固定的分量, 系数矩阵 $K_N$ 是 $2m\times2m$, 即 $4N$ 阶.
两列右端共用同一系数矩阵, 所以只须分解该矩阵一次.
从第二组方程逐指标代入第一组,
\[
\begin{aligned}
\mathbf v_j
&=f_j-\sum\limits_{k=1}^{m}A_{jk}
\left(g_k+\sum\limits_{l=1}^{m}B_{kl}\mathbf v_l\right)\\
&=f_j-\sum\limits_{k=1}^{m}A_{jk}g_k
-\sum\limits_{l=1}^{m}\left(\sum\limits_{k=1}^{m}A_{jk}B_{kl}\right)\mathbf v_l.
\end{aligned}
\]
把含未知列的项移到左边, 就是
\[
\sum\limits_{l=1}^{m}(L_N)_{jl}\mathbf v_l
=f_j-\sum\limits_{k=1}^{m}A_{jk}g_k,
\]
其中
\[
(L_N)_{jl}=\delta_{jl}+
\sum\limits_{k=1}^{m}
\dfrac{\conj{C_k}C_l}
{(\zeta_j-\conj{\zeta_k})(\conj{\zeta_k}-\zeta_l)}.
\]
这里 $\delta_{jl}$ 为 Kronecker 符号: $j=l$ 时等于 1, 否则等于 0.
这正是矩阵乘积 $L_N=I_m+AB$ 的逐元表达式.
分块消元左乘矩阵的行列式为 1, 因而
$\det K_N=\det L_N$.
所以 $\det L_N\ne0$ 同时保证两个向量分量的唯一可解性,
而不只是其中某个标量方程有解.

右端也可逐分量核对:
\[
\begin{aligned}
(\boldsymbol f-A\boldsymbol g)_{j1}
&=1+\ii q_0\sum\limits_{k=1}^{m}
\dfrac{\conj{C_k}}{(\zeta_j-\conj{\zeta_k})\conj{\zeta_k}}=r_j,\\
(\boldsymbol f-A\boldsymbol g)_{j2}
&=-\dfrac{\ii q_0}{\zeta_j}
-\sum\limits_{k=1}^{m}\dfrac{\conj{C_k}}{\zeta_j-\conj{\zeta_k}}.
\end{aligned}
\]
因此先求 $\boldsymbol v=L_N^{-1}(\boldsymbol f-A\boldsymbol g)$,
再按 $\boldsymbol u=\boldsymbol g+B\boldsymbol v$ 代回,
就得到了全部 $4N$ 条正则列, 并可直接代入上面的四个矩阵元.

\paragraph{光场的边框行列式不省略消元步骤.}
因为 $\boldsymbol v e_1=L_N^{-1}\boldsymbol r$,
\[
\sum\limits_{j=1}^{m}C_j(\mathbf v_j)_1
=\boldsymbol C^T\boldsymbol v e_1
=\boldsymbol C^T L_N^{-1}\boldsymbol r.
\]
此处是转置 $T$, 没有对 $C_j$ 再作共轭.
由 $M_{12}$ 的 $z^{-1}$ 系数乘以 $\ii$, 就得到
$q=q_0+\ii\boldsymbol C^T L_N^{-1}\boldsymbol r$.
再作一次分块消元,
\[
\begin{aligned}
&
\begin{pmatrix}I_m&0\\\ii\boldsymbol C^T L_N^{-1}&1\end{pmatrix}
\begin{pmatrix}L_N&\boldsymbol r\\-\ii\boldsymbol C^T&q_0\end{pmatrix}\\
&\qquad=\begin{pmatrix}L_N&\boldsymbol r\\0&q_0+\ii\boldsymbol C^T L_N^{-1}\boldsymbol r\end{pmatrix}.
\end{aligned}
\]
左乘矩阵行列式仍为 1. 对右边分块三角矩阵取行列式,
再除以非零的 $\det L_N$, 得到 \eqref{eq:N-determinant-q}.
这给出任意有限 $N$ 的直接代数求解, 不需要把单轨道公式反复相加.

\paragraph{换叶留数中的链式因子与负号.}
右乘 $\sigma_1$ 交换两列, 且
$\left(\dfrac{\ii z}{q_0}\sigma_1\right)^{-1}
=-\dfrac{\ii q_0}{z}\sigma_1$.
所以
\[
\widehat M_1(z)=-\dfrac{\ii q_0}{z}M_2(h(z)),\qquad
\widehat M_2(z)=-\dfrac{\ii q_0}{z}M_1(h(z)).
\]
在 $z=\zeta_j$ 附近,
\[
h(z)-\conj{\zeta_{\widehat j}}
=h'(\zeta_j)(z-\zeta_j)+O((z-\zeta_j)^2).
\]
$M_1$ 在共轭点的留数为 $-\conj{C_{\widehat j}}\mathbf u_{\widehat j}$,
与前面的 $-\ii q_0/z$ 相乘给出原证明的正号.
再用
\[
\conj{C_{\widehat j}}=-\dfrac{q_0^2}{\zeta_j^2}C_j,
\qquad h'(\zeta_j)=\dfrac{q_0^2}{\zeta_j^2},
\]
得到原证明中单行等式的每一个因子.
这一步是在完整轨道中配对, 因而对任意 $N$ 都成立.

\paragraph{行列式与分支点的正则性分开检查.}
在 $z=\zeta_j$ 处写第二列为
$\dfrac{C_jM_1(\zeta_j)}{z-\zeta_j}+O(1)$,
则行列式的留数为
\[
C_j\det(M_1(\zeta_j),M_1(\zeta_j))=0.
\]
共轭点同理. 因此行列式只可能在 $0$ 处有二阶极点,
其最高阶系数为
\[
\det\left(-\dfrac{\ii q_0}{z}\sigma_1\right)
=\left(-\dfrac{\ii q_0}{z}\right)^2\det\sigma_1
=\dfrac{q_0^2}{z^2}.
\]
这与无穷远极限 1 一起给出原证明中的 $1+\dfrac{d}{z}+\dfrac{q_0^2}{z^2}$.
换叶后的行列式比较消去 $d$, 从而 $\det M=\gamma$.
在 $z_b=\pm\ii q_0$ 处,
\[
\gamma'(z_b)=-\dfrac{2q_0^2}{z_b^3}\ne0.
\]
这里有一阶零点的是行列式, 不是说矩阵的四个元同时为零.
若四个元都在 $z_b$ 为零, 它们都含因子 $z-z_b$,
行列式便含 $(z-z_b)^2$, 与上述导数矛盾.
因此正则矩阵 $M(z_b)$ 的秩恰为一.

\paragraph{两个共振值以及一般公式的可逆性范围.}
记换叶因子仍为原文的 $\dfrac{\ii z}{q_0}\sigma_1$.
在共振点两叶上, 由 $\lambda(h(z))=-\lambda(z)$,
\[
\begin{aligned}
&M(h(z))\lambda(h(z))\sg M(h(z))^{-1}\\
&\quad=M(z)\sigma_1(-\lambda(z)\sg)\sigma_1M(z)^{-1}
=M(z)\lambda(z)\sg M(z)^{-1}.
\end{aligned}
\]
所以同一部分分式矩阵给出唯一的介质矩阵,
而 $\gamma(z_a^\pm)>0$ 保证这里所用的二阶逆矩阵存在.
计算时先完成有限系统, 再代入两个实共振原像;
无须在共振点计算含极点的传播指数.

一般 $N$ 的公式以 $\det L_N(x,t)\ne0$ 为条件.
它说明在这些点上矩阵解存在且唯一,
并不由相位乘积条件自动推出 $\det L_N$ 在全平面不为零.
下节取 $N=1$ 并指定纯虚数据后,
另行分解行列式并证明严格正性;
该全局正性证明与一般 $N$ 的条件性代数公式应分别阅读.
% END EXPANSION 08-N

上述公式将 $N$ 轨道问题化为 $2N$ 阶矩阵 $L_N$ 的求逆. 公式适用于 $\det L_N\ne0$ 的区域; 下节对一个纯虚轨道进一步显式消元, 并证明该行列式在全体实 $x,t$ 上严格为正.

\subsection{纯虚单轨道数据与全局矩阵解}

% BEGIN EXPANSION 08-E
\paragraph{单轨道的两点系统逐行消元.}
这里取 $N=1$, 因而 $m=2$. 无反射意味着 $G=0$, 因而 $W=0$.
一般表示中连续谱积分逐项等于零, 剩下的不是积分方程, 而是有限个列向量的线性代数方程.
当 $j,k=1,2$ 时, 四行逐行写为
\[
\begin{aligned}
\mathbf v_1
&=f_1- A_{11}\mathbf u_1- A_{12}\mathbf u_2,\\
\mathbf v_2
&=f_2- A_{21}\mathbf u_1- A_{22}\mathbf u_2,\\
\mathbf u_1
&=g_1+ B_{11}\mathbf v_1+ B_{12}\mathbf v_2,\\
\mathbf u_2
&=g_2+ B_{21}\mathbf v_1+ B_{22}\mathbf v_2.
\end{aligned}
\]
把后两行逐个代进第一行, 得
\[
\begin{aligned}
\mathbf v_1
={}&f_1- A_{11}g_1- A_{12}g_2\\
&-( A_{11} B_{11}
+ A_{12} B_{21})\mathbf v_1\\
&-( A_{11} B_{12}
+ A_{12} B_{22})\mathbf v_2.
\end{aligned}
\]
同样代进第二行, 得
\[
\begin{aligned}
\mathbf v_2
={}&f_2- A_{21}g_1- A_{22}g_2\\
&-( A_{21} B_{11}
+ A_{22} B_{21})\mathbf v_1\\
&-( A_{21} B_{12}
+ A_{22} B_{22})\mathbf v_2.
\end{aligned}
\]
把含未知列的项移到左侧, 就是
\[
(I_2+ A B)\boldsymbol v
=\boldsymbol f- A\boldsymbol g.
\]
这里 $I_2+ A B$ 的下标是离散点下标;
每个未知量本身仍是一条二维列.
对每个向量分量分别取值, 所以同一个二阶系数矩阵解两次.
这也解释原文将 $\boldsymbol v$ 记成两行两列数组时, 每一行实际是一条正则列的转置.
% END EXPANSION 08-E
取一个纯虚离散轨道
\begin{equation}\label{eq:imaginarypoint}
k_1=\ii b,\qquad b>q_0,\qquad
\kappa=\sqrt{b^2-q_0^2},\qquad W=b+\kappa,\qquad h_0=\dfrac{q_0}{W}\in(0,1).
\end{equation}
两个 $D_+$ 零点与透射函数为
\begin{equation}\label{eq:onezeros}
\zeta_1=\ii W,\qquad \zeta_2=-\ii q_0h_0=h(\conj{\zeta_1}),\qquad
s_{22}(z)=\dfrac{\lambda(z)-\ii\kappa}{\lambda(z)+\ii\kappa}.
\end{equation}
这正是无反射迹公式的两因子乘积, 其在 $0,\infty$ 均趋于 1. 在外部零点选择 $b_1(0)=1$, 并由换叶约化确定另一常数:
\begin{equation}\label{eq:onec}
c_1(0)=\dfrac{1}{s_{22}'(\zeta_1)}=\dfrac{4\ii\kappa}{1+h_0^2},\qquad
c_2(0)=-h_0^2c_1(0).
\end{equation}
令
\begin{equation}\label{eq:orbit-phase}
s=\kappa[x-F(\ii b)t],\qquad T=e^{2s}.
\end{equation}
由 $\lambda(\zeta_1)=\lambda(\zeta_2)=\ii\kappa$, $k(\zeta_2)=-\ii b$ 及 $F(-\ii b)=\conj{F(\ii b)}$, 得 $C_1=c_1(0)T$、$C_2=-h_0^2c_1(0)\conj{T}$.

% BEGIN EXPANSION 08-F
\paragraph{谱点, 换叶, 透射导数与四个范数系数的实算.}
由 $\kappa^2=b^2-q_0^2$ 得
\[
(b+\kappa)(b-\kappa)=q_0^2.
\]
结合 $W=b+\kappa$ 与 $h_0=\dfrac{q_0}{W}$, 逐式得到
\[
\begin{aligned}
W&=\dfrac{q_0}{h_0},&
b-\kappa&=q_0h_0,\\
2b&=\dfrac{q_0}{h_0}+q_0h_0
=\dfrac{q_0(1+h_0^2)}{h_0},\\
2\kappa&=\dfrac{q_0}{h_0}-q_0h_0
=\dfrac{q_0(1-h_0^2)}{h_0}.
\end{aligned}
\]
由于 $b>q_0>0$, 有 $\kappa>0$ 和 $W>q_0$, 所以 $0<h_0<1$.
四个点逐一列出是
\[
\zeta_1=\dfrac{\ii q_0}{h_0},\quad
\zeta_2=-\ii q_0h_0,\quad
\conj{\zeta_1}=-\dfrac{\ii q_0}{h_0},\quad
\conj{\zeta_2}=\ii q_0h_0.
\]
在第一个点代入均匀化公式:
\[
\begin{aligned}
\lambda(\zeta_1)
&=\dfrac{1}{2}\left(\dfrac{\ii q_0}{h_0}
+\dfrac{q_0^2}{\dfrac{\ii q_0}{h_0}}\right)
=\dfrac{\ii q_0}{2}(h_0^{-1}-h_0)=\ii\kappa,\\
k(\zeta_1)
&=\dfrac{\ii q_0}{2}(h_0^{-1}+h_0)=\ii b.
\end{aligned}
\]
在第二个点同样计算:
\[
\begin{aligned}
\lambda(\zeta_2)
&=\dfrac{1}{2}\left(-\ii q_0h_0
+\dfrac{q_0^2}{-\ii q_0h_0}\right)
=\dfrac{\ii q_0}{2}(h_0^{-1}-h_0)=\ii\kappa,\\
k(\zeta_2)
&=-\dfrac{\ii q_0}{2}(h_0+h_0^{-1})=-\ii b.
\end{aligned}
\]
共轭点的对应值则是
\[
\begin{aligned}
\lambda(\conj{\zeta_1})&=-\ii\kappa,&
k(\conj{\zeta_1})&=-\ii b,\\
\lambda(\conj{\zeta_2})&=-\ii\kappa,&
k(\conj{\zeta_2})&=\ii b.
\end{aligned}
\]
两个零因子的分子相乘为
\[
\begin{aligned}
(z-\zeta_1)(z-\zeta_2)
&=\left(z-\dfrac{\ii q_0}{h_0}\right)
\left(z+\ii q_0h_0\right)\\
&=z^2-\ii q_0(h_0^{-1}-h_0)z+q_0^2\\
&=z^2-2\ii\kappa z+q_0^2
=2z(\lambda-\ii\kappa).
\end{aligned}
\]
共轭分母相乘同理为 $2z(\lambda+\ii\kappa)$.
两者相除便是原文的 $s_{22}$.
对该式按 $z$ 求导, 使用商法则:
\[
\begin{aligned}
s_{22}'(z)
&=\dfrac{\lambda'(z)(\lambda+\ii\kappa)
-(\lambda-\ii\kappa)\lambda'(z)}
{(\lambda+\ii\kappa)^2}\\
&=\dfrac{2\ii\kappa\lambda'(z)}
{(\lambda+\ii\kappa)^2},\\
\lambda'(z)&=\dfrac{1}{2}\left(1-\dfrac{q_0^2}{z^2}\right).
\end{aligned}
\]
四点上的导数分别为
\[
\begin{aligned}
\lambda'(\zeta_1)=\lambda'(\conj{\zeta_1})
&=\dfrac{1+h_0^2}{2},\\
\lambda'(\zeta_2)=\lambda'(\conj{\zeta_2})
&=\dfrac{1+h_0^{-2}}{2}
=\dfrac{1+h_0^2}{2h_0^2}.
\end{aligned}
\]
所以
\[
\begin{aligned}
s_{22}'(\zeta_1)
&=\dfrac{\lambda'(\zeta_1)}{2\ii\kappa}
=\dfrac{1+h_0^2}{4\ii\kappa},\\
s_{22}'(\zeta_2)
&=\dfrac{1+h_0^2}{4\ii\kappa h_0^2}.
\end{aligned}
\]
原选择 $b_1(0)=1$ 与约化 $b_2(0)=-\conj{b_1(0)}=-1$ 给出
\[
\begin{aligned}
c_1(0)&=\dfrac{1}{s_{22}'(\zeta_1)}
=\dfrac{4\ii\kappa}{1+h_0^2},\\
c_2(0)&=\dfrac{-1}{s_{22}'(\zeta_2)}
=-\dfrac{4\ii\kappa h_0^2}{1+h_0^2}
=-h_0^2c_1(0).
\end{aligned}
\]
因此 $\conj{c_1(0)}=-c_1(0)$, $\conj{c_2(0)}=-c_2(0)$.
再计算时间指数:
\[
\begin{aligned}
2\ii\theta(\zeta_1)
&=2\ii(-\ii\kappa)[x-F(\ii b)t]
=2\kappa[x-F(\ii b)t]=2s,\\
2\ii\theta(\zeta_2)
&=2\kappa[x-F(-\ii b)t]=2\conj s.
\end{aligned}
\]
所以四个实际留数标量是
\[
\begin{aligned}
C_1&=\dfrac{4\ii\kappa}{1+h_0^2}T,&
C_2&=-\dfrac{4\ii\kappa h_0^2}{1+h_0^2}\conj T,\\
\conj{C_1}&=-\dfrac{4\ii\kappa}{1+h_0^2}\conj T,&
\conj{C_2}&=\dfrac{4\ii\kappa h_0^2}{1+h_0^2}T.
\end{aligned}
\]
% END EXPANSION 08-F
\begin{lemma}[一轨道系统的消元]\label{lem:direct-solve}
上述留数系统对任意实 $x,t$ 唯一可解. 记
\begin{equation}\label{eq:delta-pair}
\Delta_\pm=(1+h_0^2)(1+|T|^2)\pm2h_0(T+\conj{T}),
\end{equation}
则其解为
\begin{align}
\mathbf v_1&=\dfrac{1+h_0^2}{\Delta_-}
\begin{pmatrix}1-h_0\conj{T}\\\conj{T}-h_0\end{pmatrix},&
\mathbf u_1&=\dfrac{1+h_0^2}{\Delta_-}
\begin{pmatrix}h_0-T\\1-h_0T\end{pmatrix},\label{eq:direct-vu1}\\
\mathbf v_2&=\dfrac{1}{h_0}\mathbf u_1,&
\mathbf u_2&=-\dfrac{1}{h_0}\mathbf v_1.\label{eq:direct-vu2}
\end{align}
\end{lemma}
\begin{proof}
代入谱点和范数常数得到
\begin{equation}\label{eq:one-AB}
A=\dfrac{1}{1+h_0^2}
\begin{pmatrix}(h_0^2-1)\conj{T}&2h_0^2T\\-2\conj{T}&(h_0^2-1)T\end{pmatrix},\quad
B=\dfrac{1}{1+h_0^2}
\begin{pmatrix}(h_0^2-1)T&2h_0^2\conj{T}\\-2T&(h_0^2-1)\conj{T}\end{pmatrix}.
\end{equation}
对应右端为
\[
f_1=\begin{pmatrix}1\\-h_0\end{pmatrix},\quad
f_2=\begin{pmatrix}1\\h_0^{-1}\end{pmatrix},\quad
g_1=\begin{pmatrix}h_0\\1\end{pmatrix},\quad
g_2=\begin{pmatrix}-h_0^{-1}\\1\end{pmatrix}.
\]
两个向量分量分别满足同一个四阶标量系数系统,
\[
\begin{pmatrix}I_2&A\\-B&I_2\end{pmatrix}
\begin{pmatrix}\boldsymbol v\\\boldsymbol u\end{pmatrix}
=\begin{pmatrix}\boldsymbol f\\\boldsymbol g\end{pmatrix}.
\]
其 Schur 补行列式为
\begin{equation}\label{eq:one-schur-det}
\det(I_2+AB)=\dfrac{\Delta_-\Delta_+}{(1+h_0^2)^2}>0.
\end{equation}
事实上,
\[
\Delta_-=|T-h_0|^2+|1-h_0T|^2,\qquad
\Delta_+=|T+h_0|^2+|1+h_0T|^2,
\]
且 $\Delta_\pm\ge(1-h_0)^2(1+|T|^2)$. 令 $L=I_2+AB$, 用二阶逆矩阵公式计算 $L^{-1}(\boldsymbol f-A\boldsymbol g)$, 约去公因子 $\Delta_+$, 得
\[
\boldsymbol v=\dfrac{1+h_0^2}{\Delta_-}
\begin{pmatrix}1-h_0\conj{T}&\conj{T}-h_0\\
\dfrac{h_0-T}{h_0}&\dfrac{1-h_0T}{h_0}\end{pmatrix}.
\]
此处每一行是一个正则列的转置. 由 $\boldsymbol u=\boldsymbol g+B\boldsymbol v$ 得到其余两列, 即 \eqref{eq:direct-vu1}--\eqref{eq:direct-vu2}.
\end{proof}

% BEGIN EXPANSION 08-G
\subsubsection*{留数系统的八个系数, Schur 补的四元与二阶逆矩阵}
四种差分母逐个为
\[
\begin{aligned}
\zeta_1-\conj{\zeta_1}&=\dfrac{2\ii q_0}{h_0},\\
\zeta_1-\conj{\zeta_2}&=\dfrac{\ii q_0(1-h_0^2)}{h_0},\\
\zeta_2-\conj{\zeta_1}&=\dfrac{\ii q_0(1-h_0^2)}{h_0},\\
\zeta_2-\conj{\zeta_2}&=-2\ii q_0h_0.
\end{aligned}
\]
由 $2\kappa h_0=q_0(1-h_0^2)$, 八个商逐个化简为
\[
\begin{aligned}
 A_{11}
&=\dfrac{-4\ii\kappa\conj T}{1+h_0^2}
\dfrac{h_0}{2\ii q_0}
=\dfrac{(h_0^2-1)\conj T}{1+h_0^2},\\
 A_{12}
&=\dfrac{4\ii\kappa h_0^2T}{1+h_0^2}
\dfrac{h_0}{\ii q_0(1-h_0^2)}
=\dfrac{2h_0^2T}{1+h_0^2},\\
 A_{21}
&=\dfrac{-4\ii\kappa\conj T}{1+h_0^2}
\dfrac{h_0}{\ii q_0(1-h_0^2)}
=-\dfrac{2\conj T}{1+h_0^2},\\
 A_{22}
&=\dfrac{4\ii\kappa h_0^2T}{1+h_0^2}
\dfrac{-1}{2\ii q_0h_0}
=\dfrac{(h_0^2-1)T}{1+h_0^2},\\
 B_{11}
&=\dfrac{4\ii\kappa T}{1+h_0^2}
\dfrac{-h_0}{2\ii q_0}
=\dfrac{(h_0^2-1)T}{1+h_0^2},\\
 B_{12}
&=\dfrac{-4\ii\kappa h_0^2\conj T}{1+h_0^2}
\dfrac{-h_0}{\ii q_0(1-h_0^2)}
=\dfrac{2h_0^2\conj T}{1+h_0^2},\\
 B_{21}
&=\dfrac{4\ii\kappa T}{1+h_0^2}
\dfrac{-h_0}{\ii q_0(1-h_0^2)}
=-\dfrac{2T}{1+h_0^2},\\
 B_{22}
&=\dfrac{-4\ii\kappa h_0^2\conj T}{1+h_0^2}
\dfrac{1}{2\ii q_0h_0}
=\dfrac{(h_0^2-1)\conj T}{1+h_0^2}.
\end{aligned}
\]
原文已经定义 $L=I_2+ A B$.
矩阵乘法每个元只包含两个乘积, 全部写为
\[
\begin{aligned}
L_{11}
&=1+ A_{11} B_{11}
+ A_{12} B_{21}\\
&=1+\dfrac{(h_0^2-1)^2|T|^2-4h_0^2T^2}
{(1+h_0^2)^2},\\
L_{12}
&= A_{11} B_{12}
+ A_{12} B_{22}\\
&=\dfrac{2h_0^2(h_0^2-1)((\conj{T})^2+|T|^2)}
{(1+h_0^2)^2},\\
L_{21}
&= A_{21} B_{11}
+ A_{22} B_{21}\\
&=-\dfrac{2(h_0^2-1)(|T|^2+T^2)}
{(1+h_0^2)^2},\\
L_{22}
&=1+ A_{21} B_{12}
+ A_{22} B_{22}\\
&=1+\dfrac{(h_0^2-1)^2|T|^2-4h_0^2(\conj{T})^2}
{(1+h_0^2)^2}.
\end{aligned}
\]
先分别计算两个系数矩阵的行列式:
\[
\begin{aligned}
\det A
&=\dfrac{(h_0^2-1)^2\conj T T+4h_0^2T\conj T}
{(1+h_0^2)^2}\\
&=\dfrac{(h_0^4-2h_0^2+1+4h_0^2)|T|^2}
{(1+h_0^2)^2}=|T|^2,\\
\det B&=|T|^2.
\end{aligned}
\]
任意二阶矩阵的恒等式
$\det(I_2+ A B)
=1+\tr( A B)+\det( A B)$
可由 $L_{11}L_{22}-L_{12}L_{21}$ 直接展开得到.
这里代入上面的四元, 给出
\[
\det L
=1+|T|^4
+\dfrac{2(h_0^2-1)^2|T|^2
-4h_0^2(T^2+(\conj{T})^2)}{(1+h_0^2)^2}.
\]
另一方面, 把 $\Delta_\pm$ 的乘积按平方差展开:
\[
\begin{aligned}
\Delta_-\Delta_+
={}&(1+h_0^2)^2(1+|T|^2)^2
-4h_0^2(T+\conj T)^2\\
={}&(1+h_0^2)^2(1+|T|^4)\\
&+\bigl(2(1+h_0^2)^2-8h_0^2\bigr)|T|^2
-4h_0^2(T^2+(\conj{T})^2)\\
={}&(1+h_0^2)^2(1+|T|^4)
+2(h_0^2-1)^2|T|^2
-4h_0^2(T^2+(\conj{T})^2).
\end{aligned}
\]
除以 $(1+h_0^2)^2$ 恰为 $\det L$.
正性可不用任何关于 $T$ 相位的限制:
\[
\begin{aligned}
|T\mp h_0|^2+|1\mp h_0T|^2
&=(1+h_0^2)(1+|T|^2)
\mp2h_0(T+\conj T),\\
\Delta_\pm
&\ge(1+h_0^2)(1+|T|^2)-4h_0|T|\\
&=(1-h_0)^2(1+|T|^2)
+2h_0(|T|-1)^2\\
&\ge(1-h_0)^2(1+|T|^2)>0.
\end{aligned}
\]
因此四阶系数系统对于任意复数 $T$ 都可逆;
实际 $T=e^{2s}$ 当然也在这个范围内.

\paragraph{右端四个标量以及伴随矩阵乘法.}
将两条列的转置排成行, 原有右端的四个元分别为
\[
\begin{aligned}
(\boldsymbol f- A\boldsymbol g)_{11}
&=1-h_0 A_{11}+h_0^{-1} A_{12}\\
&=\dfrac{1+h_0^2-h_0(h_0^2-1)\conj T+2h_0T}
{1+h_0^2},\\
(\boldsymbol f- A\boldsymbol g)_{12}
&=-h_0- A_{11}- A_{12}\\
&=\dfrac{-h_0(1+h_0^2)-(h_0^2-1)\conj T-2h_0^2T}
{1+h_0^2},\\
(\boldsymbol f- A\boldsymbol g)_{21}
&=1-h_0 A_{21}+h_0^{-1} A_{22}\\
&=\dfrac{h_0(1+h_0^2)+2h_0^2\conj T+(h_0^2-1)T}
{h_0(1+h_0^2)},\\
(\boldsymbol f- A\boldsymbol g)_{22}
&=h_0^{-1}- A_{21}- A_{22}\\
&=\dfrac{1+h_0^2+2h_0\conj T-h_0(h_0^2-1)T}
{h_0(1+h_0^2)}.
\end{aligned}
\]
二阶逆矩阵的完整式为
\[
L^{-1}
=\dfrac{1}{\det L}
\begin{pmatrix}L_{22}&-L_{12}\\-L_{21}&L_{11}\end{pmatrix}.
\]
例如其与第一列相乘的第一元, 在通分后是
\[
\begin{aligned}
&L_{22}(\boldsymbol f- A\boldsymbol g)_{11}
-L_{12}(\boldsymbol f- A\boldsymbol g)_{21}\\
={}&\dfrac{1}{(1+h_0^2)^3}
\Bigl[
\bigl((1+h_0^2)^2+(h_0^2-1)^2|T|^2-4h_0^2(\conj{T})^2\bigr)\\
&\qquad\cdot
\bigl(1+h_0^2-h_0(h_0^2-1)\conj T+2h_0T\bigr)\\
&\qquad-2h_0(h_0^2-1)((\conj{T})^2+|T|^2)
\bigl(h_0(1+h_0^2)+2h_0^2\conj T+(h_0^2-1)T\bigr)
\Bigr]\\
={}&\dfrac{\Delta_+}{1+h_0^2}(1-h_0\conj T).
\end{aligned}
\]
最后一行可以逐幂核验: 将 $|T|^2$ 写回 $T\conj T$,
同类项相加, 分子正好为
\[
(1+h_0^2)^2(1-h_0\conj T)
\bigl((1+h_0^2)(1+T\conj T)+2h_0(T+\conj T)\bigr).
\]
同一伴随矩阵乘法的其余三个元, 不省略其组合方式:
\[
\begin{aligned}
&L_{22}(\boldsymbol f- A\boldsymbol g)_{12}
-L_{12}(\boldsymbol f- A\boldsymbol g)_{22}
=\dfrac{\Delta_+}{1+h_0^2}(\conj T-h_0),\\
&-L_{21}(\boldsymbol f- A\boldsymbol g)_{11}
+L_{11}(\boldsymbol f- A\boldsymbol g)_{21}
=\dfrac{\Delta_+}{h_0(1+h_0^2)}(h_0-T),\\
&-L_{21}(\boldsymbol f- A\boldsymbol g)_{12}
+L_{11}(\boldsymbol f- A\boldsymbol g)_{22}
=\dfrac{\Delta_+}{h_0(1+h_0^2)}(1-h_0T).
\end{aligned}
\]
用 $\det L=\dfrac{\Delta_-\Delta_+}{(1+h_0^2)^2}$ 除这四项,
每一项中的 $\Delta_+$ 都约掉, 正好得到原文的 $\boldsymbol v$.

\paragraph{将 $\boldsymbol v$ 代回以得到 $\boldsymbol u$.}
例如第一条列的第一元是
\[
\begin{aligned}
(\mathbf u_1)_1
&=h_0+ B_{11}(\mathbf v_1)_1
+ B_{12}(\mathbf v_2)_1\\
&=h_0+\dfrac{(h_0^2-1)T(1-h_0\conj T)
+2h_0\conj T(h_0-T)}{\Delta_-}\\
&=\dfrac{h_0\Delta_-+(h_0^2-1)T(1-h_0\conj T)
+2h_0\conj T(h_0-T)}{\Delta_-}\\
&=\dfrac{(1+h_0^2)(h_0-T)}{\Delta_-}.
\end{aligned}
\]
其余三个元使用同一个代入, 分子分别为
\[
\begin{aligned}
\Delta_-(\mathbf u_1)_2
&=\Delta_-+(h_0^2-1)T(\conj T-h_0)
+2h_0\conj T(1-h_0T)\\
&=(1+h_0^2)(1-h_0T),\\
\Delta_-(\mathbf u_2)_1
&=-\dfrac{\Delta_-}{h_0}-2T(1-h_0\conj T)
+\dfrac{h_0^2-1}{h_0}\conj T(h_0-T)\\
&=-\dfrac{1+h_0^2}{h_0}(1-h_0\conj T),\\
\Delta_-(\mathbf u_2)_2
&=\Delta_--2T(\conj T-h_0)
+\dfrac{h_0^2-1}{h_0}\conj T(1-h_0T)\\
&=-\dfrac{1+h_0^2}{h_0}(\conj T-h_0).
\end{aligned}
\]
这四式既给出 $\mathbf u_1,\mathbf u_2$, 也直接检查
$\mathbf v_2=\dfrac{\mathbf u_1}{h_0}$ 与 $\mathbf u_2=-\dfrac{\mathbf v_1}{h_0}$.
没有使用预先假定的换叶关系去缩小系统;
换叶关系是解出完整系统后得到的结论.
% END EXPANSION 08-G
将所得列向量代入部分分式表示. 其系数可用以下秩一正交投影写成紧凑形式:
\begin{equation}\label{eq:algebraic-P}
P=\dfrac{1}{\Delta_-}
\begin{pmatrix}
|T-h_0|^2&(T-h_0)(1-h_0\conj{T})\\
(1-h_0T)(\conj{T}-h_0)&|1-h_0T|^2
\end{pmatrix},\qquad H=\diag(\ii,-\ii).
\end{equation}
由 $\Delta_-=|T-h_0|^2+|1-h_0T|^2$, 有 $P^\dagger=P$, $P^2=P$, $\tr P=1$.

\begin{theorem}[完整四点轨道的矩阵解]\label{thm:oneRHP}\label{prop:unique-reflectionless}
在背景参数 $q_0>0$, $a\ne0$, $\eta_0\ne0$, $\alpha,\beta\in\mathbb R$ 和 $b>q_0$ 下, 数据 \eqref{eq:onezeros}--\eqref{eq:onec} 的无反射 RHP 有唯一解
\begin{equation}\label{eq:explicitM}
M(z)=H\bigl[(k(z)+\ii b)I-2\ii bP\bigr]Y(z)
\diag\left(\dfrac{1}{\lambda(z)+\ii\kappa},\dfrac{1}{\lambda(z)-\ii\kappa}\right)H^{-1}.
\end{equation}
它满足定义 \ref{def:RHP} 的全部条件, 在分支点、轮廓交点与共振原像解析, 在分支点秩为一, 在共振原像可逆. 唯一性取于除 $0$ 和指定四点外没有有限奇点、非极点轮廓处可去并满足留数及两处归一化的亚纯矩阵类.
\end{theorem}
\begin{proof}
在部分分式的第二列合并分母,
$(z-\zeta_1)(z-\zeta_2)=2z(\lambda-\ii\kappa)$, 第一列同理. 利用
$2bh_0=q_0(1+h_0^2)$、$2\kappa h_0=q_0(1-h_0^2)$ 化简分子, 得到 \eqref{eq:explicitM}. 留数、两处归一化及恒等跳跃已由上述留数系统 保证.

共轭矩阵 $\sigma_2\conj{M(\conj{z})}\sigma_2$ 满足同一留数系统, 其中交换两组极点产生 \eqref{eq:resminus} 的负号. 换叶矩阵 $M(h(z))C(z)^{-1}$ 也满足该系统: 留数变换的导数因子由 \eqref{eq:cpair} 抵消, 而 $\theta(h(\conj{\zeta}))=-\conj{\theta(\zeta)}$ 保证指数一致. 两处归一化相互交换后仍为规定主项. 引理 \ref{lem:direct-solve} 的唯一性遂给出 \eqref{eq:Msym}.

记
\begin{equation}\label{eq:Nfactor}
N(k)=(k+\ii b)I-2\ii bP.
\end{equation}
$P$ 的特征值为 $1,0$, 故
\[
\det N=(k-\ii b)(k+\ii b)=\lambda^2+\kappa^2.
\]
它抵消两个列分母的乘积, 得 $\det M=\det Y=\gamma$. 在分支点 $z=\pm\ii q_0$ 上, $\lambda=0$, 两个列分母均非零; $N$ 可逆而 $Y$ 秩一, 因此 $M$ 解析且秩一. 交点和共振原像的 $\lambda$ 为实数, 也没有列分母零点; 共振处 $\gamma>0$, 所以 $M$ 可逆. 显式表达式不再含有其他有限奇点.

任一属于所述函数类的解均具有 \eqref{eq:partialfraction} 的表示, 因而由同一留数系统确定. Schur 行列式 \eqref{eq:one-schur-det} 严格为正, 保证该解唯一.
\end{proof}

% BEGIN EXPANSION 08-H
\subsubsection*{投影, 有理矩阵的四元以及四点留数的直接核验}
先逐项检查原有投影矩阵.
其四元是
\[
\begin{aligned}
P_{11}&=\dfrac{(T-h_0)(\conj T-h_0)}{\Delta_-},\\
P_{12}&=\dfrac{(T-h_0)(1-h_0\conj T)}{\Delta_-},\\
P_{21}&=\dfrac{(1-h_0T)(\conj T-h_0)}{\Delta_-},\\
P_{22}&=\dfrac{(1-h_0T)(1-h_0\conj T)}{\Delta_-}.
\end{aligned}
\]
分母为正实数, 所以 $P_{11},P_{22}$ 为实数且
$\conj{P_{12}}=P_{21}$.
这逐元给出 $P^\dagger=P$.
迹的计算为
\[
P_{11}+P_{22}
=\dfrac{|T-h_0|^2+|1-h_0T|^2}{\Delta_-}=1.
\]
行列式分子中两项完全相同:
\[
\begin{aligned}
\Delta_-^2\det P
={}&(T-h_0)(\conj T-h_0)(1-h_0T)(1-h_0\conj T)\\
&-(T-h_0)(1-h_0\conj T)(1-h_0T)(\conj T-h_0)=0.
\end{aligned}
\]
平方矩阵四元的乘法则是
\[
\begin{aligned}
(P^2)_{11}
&=\dfrac{|T-h_0|^4+|T-h_0|^2|1-h_0T|^2}{\Delta_-^2}
=P_{11},\\
(P^2)_{12}
&=\dfrac{(T-h_0)(1-h_0\conj T)
(|T-h_0|^2+|1-h_0T|^2)}{\Delta_-^2}
=P_{12},\\
(P^2)_{21}
&=\dfrac{(1-h_0T)(\conj T-h_0)
(|T-h_0|^2+|1-h_0T|^2)}{\Delta_-^2}
=P_{21},\\
(P^2)_{22}
&=\dfrac{|1-h_0T|^2|T-h_0|^2+|1-h_0T|^4}{\Delta_-^2}
=P_{22}.
\end{aligned}
\]
其中第二、第三行分子中紧邻括号表示乘积.
因此该矩阵确为正交投影.
其迹为 1 而非零, 行列式为零, 所以秩恰为一.

\paragraph{由部分分式合并得到显式矩阵, 不跳过四个元.}
四个留数列的系数分别是
\[
\begin{aligned}
C_1\mathbf v_1
&=\dfrac{4\ii\kappa T}{\Delta_-}
\begin{pmatrix}1-h_0\conj T\\\conj T-h_0\end{pmatrix},\\
C_2\mathbf v_2
&=\dfrac{4\ii\kappa h_0\conj T}{\Delta_-}
\begin{pmatrix}T-h_0\\h_0T-1\end{pmatrix},\\
-\conj{C_1}\mathbf u_1
&=\dfrac{4\ii\kappa\conj T}{\Delta_-}
\begin{pmatrix}h_0-T\\1-h_0T\end{pmatrix},\\
-\conj{C_2}\mathbf u_2
&=\dfrac{4\ii\kappa h_0T}{\Delta_-}
\begin{pmatrix}1-h_0\conj T\\\conj T-h_0\end{pmatrix}.
\end{aligned}
\]
因此原部分分式四元完整写为
\[
\begin{aligned}
M_{11}
&=1+\dfrac{4\ii\kappa}{\Delta_-}
\left[
\dfrac{\conj T(h_0-T)}{z+\dfrac{\ii q_0}{h_0}}
+\dfrac{h_0T(1-h_0\conj T)}{z-\ii q_0h_0}
\right],\\
M_{21}
&=-\dfrac{\ii q_0}{z}
+\dfrac{4\ii\kappa}{\Delta_-}
\left[
\dfrac{\conj T(1-h_0T)}{z+\dfrac{\ii q_0}{h_0}}
+\dfrac{h_0T(\conj T-h_0)}{z-\ii q_0h_0}
\right],\\
M_{12}
&=-\dfrac{\ii q_0}{z}
+\dfrac{4\ii\kappa}{\Delta_-}
\left[
\dfrac{T(1-h_0\conj T)}{z-\dfrac{\ii q_0}{h_0}}
+\dfrac{h_0\conj T(T-h_0)}{z+\ii q_0h_0}
\right],\\
M_{22}
&=1+\dfrac{4\ii\kappa}{\Delta_-}
\left[
\dfrac{T(\conj T-h_0)}{z-\dfrac{\ii q_0}{h_0}}
-\dfrac{h_0\conj T(1-h_0T)}{z+\ii q_0h_0}
\right].
\end{aligned}
\]
第一列两个分母的乘积是 $2z(\lambda+\ii\kappa)$,
第二列两个分母的乘积是 $2z(\lambda-\ii\kappa)$.
例如 $M_{12}$ 中两个留数项通分时的分子为
\[
\begin{aligned}
&T(1-h_0\conj T)(z+\ii q_0h_0)
+h_0\conj T(T-h_0)(z-\dfrac{\ii q_0}{h_0})\\
={}&z\bigl(T-h_0^2\conj T\bigr)
+\ii q_0\bigl(h_0(T+\conj T)-(1+h_0^2)|T|^2\bigr).
\end{aligned}
\]
其余三个元的对应分子也逐个为
\[
\begin{aligned}
&\conj T(h_0-T)(z-\ii q_0h_0)
+h_0T(1-h_0\conj T)(z+\dfrac{\ii q_0}{h_0})\\
&\qquad=z\bigl(h_0(T+\conj T)-(1+h_0^2)|T|^2\bigr)
+\ii q_0(T-h_0^2\conj T),\\
&\conj T(1-h_0T)(z-\ii q_0h_0)
+h_0T(\conj T-h_0)(z+\dfrac{\ii q_0}{h_0})\\
&\qquad=z(\conj T-h_0^2T)
+\ii q_0\bigl((1+h_0^2)|T|^2-h_0(T+\conj T)\bigr),\\
&T(\conj T-h_0)(z+\ii q_0h_0)
-h_0\conj T(1-h_0T)(z-\dfrac{\ii q_0}{h_0})\\
&\qquad=z\bigl((1+h_0^2)|T|^2-h_0(T+\conj T)\bigr)
+\ii q_0(\conj T-h_0^2T).
\end{aligned}
\]
接着将 $N$ 的四元直接写出:
\[
\begin{aligned}
N_{11}&=k+\ii b-2\ii bP_{11},&
N_{12}&=-2\ii bP_{12},\\
N_{21}&=-2\ii bP_{21},&
N_{22}&=k+\ii b-2\ii bP_{22}.
\end{aligned}
\]
因为
\[
H=\begin{pmatrix}\ii&0\\0&-\ii\end{pmatrix},\qquad
H^{-1}=\begin{pmatrix}-\ii&0\\0&\ii\end{pmatrix},
\]
对任意矩阵左右乘 $H,H^{-1}$ 会保持对角元, 将两个非对角元变号.
因此显式矩阵 \eqref{eq:explicitM} 按元展开为
\[
\begin{aligned}
M_{11}
&=\dfrac{N_{11}-\dfrac{\ii q_0}{z}N_{12}}{\lambda+\ii\kappa},\\
M_{21}
&=-\dfrac{N_{21}-\dfrac{\ii q_0}{z}N_{22}}{\lambda+\ii\kappa},\\
M_{12}
&=-\dfrac{N_{12}-\dfrac{\ii q_0}{z}N_{11}}{\lambda-\ii\kappa},\\
M_{22}
&=\dfrac{N_{22}-\dfrac{\ii q_0}{z}N_{21}}{\lambda-\ii\kappa}.
\end{aligned}
\]
代入 $P_{jk}$ 后, 这四式与上面的通分式逐元相同.
核对时仅需使用
\[
\lambda-k=\dfrac{q_0^2}{z},\qquad
2bh_0=q_0(1+h_0^2),\qquad
2\kappa h_0=q_0(1-h_0^2).
\]
例如第三式的分子可无省略地写为
\[
\begin{aligned}
&2\ii b(T-h_0)(1-h_0\conj T)
+\dfrac{\ii q_0}{z}
\bigl((k+\ii b)\Delta_--2\ii b|T-h_0|^2\bigr)\\
={}&-\ii q_0(\lambda-\ii\kappa)\Delta_-\\
&+\dfrac{2\ii\kappa}{z}
\left[
z(T-h_0^2\conj T)
+\ii q_0\bigl(h_0(T+\conj T)-(1+h_0^2)|T|^2\bigr)
\right].
\end{aligned}
\]
两侧同除以 $(\lambda-\ii\kappa)\Delta_-$ 就得到先前的 $M_{12}$.
这给出了从部分分式到紧凑矩阵乘积的可直接复算过程.

\paragraph{外部 $D_+$ 点 $\zeta_1$ 的正则列与留数.}
在此点 $k=\ii b$, 所以 $N=2\ii b(I-P)$.
两列背景向量分别为 $(1,-h_0)^T,(-h_0,1)^T$.
投影的两次乘法逐项为
\[
\begin{aligned}
(I-P)\begin{pmatrix}1\\-h_0\end{pmatrix}
&=\dfrac{1-h_0^2}{\Delta_-}
\begin{pmatrix}1-h_0\conj T\\-(\conj T-h_0)\end{pmatrix},\\
(I-P)\begin{pmatrix}-h_0\\1\end{pmatrix}
&=\dfrac{(1-h_0^2)T}{\Delta_-}
\begin{pmatrix}-(1-h_0\conj T)\\\conj T-h_0\end{pmatrix}.
\end{aligned}
\]
例如第一式第一元的分子是
\[
|1-h_0T|^2+h_0(T-h_0)(1-h_0\conj T)
=(1-h_0\conj T)(1-h_0^2).
\]
先读正则第一列, 分母是 $2\ii\kappa$:
\[
M_1(\zeta_1)
=\dfrac{b(1-h_0^2)}{\kappa\Delta_-}
\begin{pmatrix}1-h_0\conj T\\\conj T-h_0\end{pmatrix}
=\mathbf v_1.
\]
这里用 $\dfrac{b(1-h_0^2)}{\kappa}=1+h_0^2$.
第二列分母在此点的导数是 $\dfrac{1+h_0^2}{2}$.
注意 $H$ 对非对角元的变号后, 其留数为
\[
\begin{aligned}
\Res_{z=\zeta_1}M_2
&=\dfrac{4\ii b(1-h_0^2)T}{(1+h_0^2)\Delta_-}
\begin{pmatrix}1-h_0\conj T\\\conj T-h_0\end{pmatrix}\\
&=\dfrac{4\ii\kappa T}{\Delta_-}
\begin{pmatrix}1-h_0\conj T\\\conj T-h_0\end{pmatrix}
=C_1\mathbf v_1.
\end{aligned}
\]

\paragraph{内部 $D_+$ 点 $\zeta_2$ 的正则列与留数.}
这里 $k=-\ii b$, 所以 $N=-2\ii bP$.
两列背景向量为 $(1,h_0^{-1})^T,(h_0^{-1},1)^T$.
所需投影乘法分别是
\[
\begin{aligned}
P\begin{pmatrix}1\\h_0^{-1}\end{pmatrix}
&=\dfrac{1-h_0^2}{h_0\Delta_-}
\begin{pmatrix}T-h_0\\1-h_0T\end{pmatrix},\\
P\begin{pmatrix}h_0^{-1}\\1\end{pmatrix}
&=\dfrac{(1-h_0^2)\conj T}{h_0\Delta_-}
\begin{pmatrix}T-h_0\\1-h_0T\end{pmatrix}.
\end{aligned}
\]
两式的标量来源可逐个查看:
\[
\begin{aligned}
(\conj T-h_0)+h_0^{-1}(1-h_0\conj T)
&=\dfrac{1-h_0^2}{h_0},\\
h_0^{-1}(\conj T-h_0)+(1-h_0\conj T)
&=\dfrac{(1-h_0^2)\conj T}{h_0}.
\end{aligned}
\]
正则第一列因此是
\[
M_1(\zeta_2)
=\dfrac{1+h_0^2}{h_0\Delta_-}
\begin{pmatrix}h_0-T\\1-h_0T\end{pmatrix}
=\mathbf v_2.
\]
第二列分母导数为 $\dfrac{1+h_0^2}{2h_0^2}$, 所以
\[
\begin{aligned}
\Res_{z=\zeta_2}M_2
&=\dfrac{4\ii b h_0(1-h_0^2)\conj T}
{(1+h_0^2)\Delta_-}
\begin{pmatrix}T-h_0\\-(1-h_0T)\end{pmatrix}\\
&=\dfrac{4\ii\kappa h_0\conj T}{\Delta_-}
\begin{pmatrix}T-h_0\\h_0T-1\end{pmatrix}
=C_2\mathbf v_2.
\end{aligned}
\]

\paragraph{外部 $D_-$ 点 $\conj{\zeta_1}$ 的正则列与留数.}
这里仍有 $N=-2\ii bP$, 但 $\lambda=-\ii\kappa$.
背景两列为 $(1,h_0)^T,(h_0,1)^T$.
直接相乘得到
\[
\begin{aligned}
P\begin{pmatrix}1\\h_0\end{pmatrix}
&=\dfrac{(1-h_0^2)\conj T}{\Delta_-}
\begin{pmatrix}T-h_0\\1-h_0T\end{pmatrix},\\
P\begin{pmatrix}h_0\\1\end{pmatrix}
&=\dfrac{1-h_0^2}{\Delta_-}
\begin{pmatrix}T-h_0\\1-h_0T\end{pmatrix}.
\end{aligned}
\]
正则第二列分母为 $-2\ii\kappa$, 因而
\[
M_2(\conj{\zeta_1})
=\dfrac{1+h_0^2}{\Delta_-}
\begin{pmatrix}h_0-T\\1-h_0T\end{pmatrix}
=\mathbf u_1.
\]
第一列分母导数为 $\dfrac{1+h_0^2}{2}$, 所以
\[
\begin{aligned}
\Res_{z=\conj{\zeta_1}}M_1
&=\dfrac{4\ii\kappa\conj T}{\Delta_-}
\begin{pmatrix}h_0-T\\1-h_0T\end{pmatrix}\\
&=-\conj{C_1}\mathbf u_1.
\end{aligned}
\]
其负号已经由 $\conj{C_1}=-\dfrac{4\ii\kappa\conj T}{1+h_0^2}$ 明确给出.

\paragraph{内部 $D_-$ 点 $\conj{\zeta_2}$ 的正则列与留数.}
此时 $N=2\ii b(I-P)$, $\lambda=-\ii\kappa$,
背景两列为 $(1,-h_0^{-1})^T,(-h_0^{-1},1)^T$.
直接计算
\[
\begin{aligned}
(I-P)\begin{pmatrix}1\\-h_0^{-1}\end{pmatrix}
&=\dfrac{(1-h_0^2)T}{h_0\Delta_-}
\begin{pmatrix}1-h_0\conj T\\-(\conj T-h_0)\end{pmatrix},\\
(I-P)\begin{pmatrix}-h_0^{-1}\\1\end{pmatrix}
&=\dfrac{1-h_0^2}{h_0\Delta_-}
\begin{pmatrix}-(1-h_0\conj T)\\\conj T-h_0\end{pmatrix}.
\end{aligned}
\]
正则第二列是
\[
M_2(\conj{\zeta_2})
=-\dfrac{1+h_0^2}{h_0\Delta_-}
\begin{pmatrix}1-h_0\conj T\\\conj T-h_0\end{pmatrix}
=\mathbf u_2.
\]
第一列分母导数为 $\dfrac{1+h_0^2}{2h_0^2}$, 给出
\[
\begin{aligned}
\Res_{z=\conj{\zeta_2}}M_1
&=\dfrac{4\ii\kappa h_0T}{\Delta_-}
\begin{pmatrix}1-h_0\conj T\\\conj T-h_0\end{pmatrix}\\
&=-\conj{C_2}\mathbf u_2.
\end{aligned}
\]
至此四个点全部按显式矩阵直接取过正则列和留数,
不再只依赖部分分式的抽象结论.

\paragraph{$N$ 的投影分解, 逆矩阵与行列式.}
利用 $I=P+(I-P)$, 原有 $N$ 可分解为
\[
N=(k-\ii b)P+(k+\ii b)(I-P).
\]
因为 $P(I-P)=(I-P)P=0$, 直接相乘得到
\[
\begin{aligned}
N^{-1}
&=\dfrac{P}{k-\ii b}+\dfrac{I-P}{k+\ii b}\\
&=\dfrac{(k+\ii b)P+(k-\ii b)(I-P)}{k^2+b^2}\\
&=\dfrac{(k-\ii b)I+2\ii bP}{k^2+b^2}.
\end{aligned}
\]
分母不为零处该式逐项成立, 其余处按亚纯函数理解.
两个互补秩一子空间上 $N$ 的特征值是 $k-\ii b,k+\ii b$,
所以
\[
\det N=k^2+b^2
=\lambda^2-q_0^2+b^2
=\lambda^2+\kappa^2.
\]
于是显式矩阵全部因子的行列式相乘为
\[
\begin{aligned}
\det M
&=\det H\,\det N\,\det Y
\dfrac{1}{(\lambda+\ii\kappa)(\lambda-\ii\kappa)}\det H^{-1}\\
&=\dfrac{(\lambda^2+\kappa^2)\gamma}{\lambda^2+\kappa^2}
=\gamma.
\end{aligned}
\]
在离散点, 该式是约去公因子后的亚纯恒等式.

\paragraph{两域约化为何由同一个有限系统固定.}
共轭变换对任意二阶矩阵的作用是
\[
\sigma_2\conj M\sigma_2
=\begin{pmatrix}
\conj{M_{22}}&-\conj{M_{21}}\\
-\conj{M_{12}}&\conj{M_{11}}
\end{pmatrix}.
\]
因此 $D_+$ 的第二列极点转成 $D_-$ 的第一列极点,
列交换时的负号正是负留数条件中的负号.
换叶关系的列形式是
\[
M_1(h(z))=\dfrac{\ii z}{q_0}M_2(z),\qquad
M_2(h(z))=\dfrac{\ii z}{q_0}M_1(z).
\]
例如 $h(\conj{\zeta_1})=\zeta_2$ 时,
$\dfrac{\ii\conj{\zeta_1}}{q_0}=h_0^{-1}$,
第一式正给出 $\mathbf v_2=\dfrac{\mathbf u_1}{h_0}$.
在 $h(\zeta_1)=\conj{\zeta_2}$ 时,
$\dfrac{\ii\zeta_1}{q_0}=-h_0^{-1}$,
第二式给出 $\mathbf u_2=-\dfrac{\mathbf v_1}{h_0}$.
这两个正则列关系已经由完整消元独立得到.
在留数中还要乘 $h'(z)=\dfrac{q_0^2}{z^2}$;
原有 $C$ 的配对公式恰好提供这个导数因子,
所以变换后的矩阵满足同一留数系统.
由于该系统可逆, 两域约化得以全局成立.

\paragraph{分支点与轮廓点的解析性.}
在 $z=\pm\ii q_0$ 处 $\lambda=0$,
$\lambda\pm\ii\kappa$ 都非零, $\det N=\kappa^2>0$.
只有 $Y$ 的行列式为零, 且
\[
Y(\ii q_0)=I-\sigma_1,\qquad
Y(-\ii q_0)=I+\sigma_1.
\]
这两个矩阵均非零且秩为一, 所以 $M$ 解析且秩为一.
在交点 $z=\pm q_0$ 处, $\lambda=\pm q_0$ 为非零实数,
列分母也不为零.
在共振原像 $z_a^\pm$ 处, $\lambda=\pm\sqrt{a^2+q_0^2}$,
列分母非零且 $\gamma>0$, 所以 $M$ 解析并可逆.
全部判断直接来自显式有理式的分母, 无需附加轮廓上的跳跃论证.
% END EXPANSION 08-H
